\documentclass[11pt]{article}
\usepackage[T1]{fontenc}
\usepackage{lmodern}
\usepackage{amsmath,amssymb,amsthm}
\usepackage[margin=1in]{geometry}
\usepackage{microtype,booktabs,longtable}
\usepackage{xurl}
\usepackage[hidelinks,breaklinks]{hyperref}
\hypersetup{pdftitle={Interval volume from thinned causal orders: finite confidence and calibration limits},pdfauthor={Daniel Eric Fredriksen}}
\newtheorem{theorem}{Theorem}[section]
\newtheorem{proposition}[theorem]{Proposition}
\newtheorem{lemma}[theorem]{Lemma}
\theoremstyle{remark}
\newtheorem{remark}[theorem]{Remark}
\newcommand{\R}{\mathbb R}
\newcommand{\K}{\mathcal K}
\newcommand{\Prb}{\mathbb P}
\newcommand{\E}{\mathbb E}
\newcommand{\TV}{\operatorname{TV}}
\newcommand{\Bin}{\operatorname{Bin}}
\newcommand{\Poi}{\operatorname{Poi}}
\newcommand{\dd}{\,\mathrm d}
\newcommand{\ind}{\mathbf 1}
\title{Interval volume from thinned causal orders:\\finite confidence and calibration limits}
\author{Daniel Eric Fredriksen\\Quantyra Inc.\\
\url{https://github.com/Quantyra/quantyra-null-cone-lean}}
\date{10 October 2026\\Version 0.1.0 --- unpublished manuscript}
\begin{document}
\maketitle
\begin{abstract}
An event detector samples a causal geometry with an unknown spatially
varying retention probability. We study the normalized physical volume of
one marked chronological interval when only the full order of the retained
events and two distinguished anchors is observed. In a smooth, bounded
class of two-dimensional conformal geometries, with known detector ratio
bound $1\le R\le2$, the fixed-confidence minimax accuracy for $n\ge1$
retained events has scale
$\min\{1,n^{-1/2}+(R-1)/(R+1)\}$, uniformly in $n$ and $R$.
The lower bound covers arbitrary measurable estimators of the full marked
order, including independent randomization. It uses admitted geometric
alternatives with either close sampling laws or exactly equal detected
laws. We separately give a finite confidence report by composing certified
binomial endpoints with the established bounded-selection transforms of
Aronow and Lee. Complete thinning arguments connect fixed, independently
mixed, Poisson and first-$n$ retention experiments. The mathematical chain
has exact-source Lean verification on Google Cloud. A preserved pilot and
a separate frozen finite comparison support retaining the exact-binomial
report for practical reporting; the minimax upper bound is conservative.
No absolute volume scale or physical detector calibration is inferred.
\end{abstract}

\section{Question and contribution}
Counting the events between two marked events is an intrinsic operation on
a causal order. Its physical interpretation depends on how events are
sampled. If detection favors some locations, the proportion of retained
events in an interval estimates its \emph{detected} probability, rather
than its physical volume fraction. More retained events alone need not
remove this discrepancy.

This paper connects three results in one explicit experiment: the law of
independent thinning, an executable finite confidence report under bounded
unknown detection, and a uniform upper and lower accuracy theorem for a
geometric class. The lower construction preserves the entire marked-order
law, not just the interval count. Consequently, an estimator cannot avoid
the calibration obstruction by using additional order relations or an
independent random seed. A second construction gives the sampling term
even under perfect detection.

The statistical ingredients have established antecedents. The binary
selection transforms are those of Aronow and Lee \cite{aronow}; confidence
inversion follows Clopper and Pearson \cite{clopper}, and the analytic
comparison uses Hoeffding's inequality \cite{hoeffding}. Partial
identification distinguishes an unknown target from bounds on its possible
values, and confidence for either object requires an additional sampling
argument \cite{imbens,stoye,tudball}. Our contribution is the connected
geometric experiment, its admitted full-law alternatives, the finite
uniform rate, and the formal and computational evidence. We make no
exhaustive priority claim for all related statistical problems.

The physical question here is narrower than reconstructing a complete
metric from order. The target is one normalized interval volume in a
specified two-dimensional model. Coordinates define the model but are not
observed. The fixed anchors and their roles are part of the experiment;
choosing an interval after looking at the sample would require a different
coverage argument.

\section{Geometry, observations and the main result}
Let $D=[0,1]^2$ and write $x=(u,v)$. A function $\rho:\R^2\to\R$
belongs to $\K$ when it is smooth on some open neighborhood of $D$ and,
for all points of $D$, satisfies
\begin{align}
 \tfrac12\le\rho\le\tfrac32,\qquad
 |\rho(x)-\rho(y)|\le2\|x-y\|_2,\label{eq:class}\\
 \int_0^1\rho(u,t)\dd t=1,\qquad
 \int_0^1\rho(t,v)\dd t=1.\nonumber
\end{align}
The Lipschitz norm is Euclidean. Only behavior on $D$ and smoothness near
it are required; positivity on all of $\R^2$ is not assumed. Set
$\mu_\rho(B)=\int_{B\cap D}\rho(u,v)\dd u\dd v$.

\begin{proposition}[Geometric interpretation]\label{prop:geometry}
The measure $\mu_\rho$ is a probability measure. On $D$, equip the model
with the Lorentzian metric
\[
 g_\rho=-\rho(\dd u\otimes\dd v+\dd v\otimes\dd u)
\]
and future direction increasing both coordinates. Its volume element is
$\rho\dd u\dd v$ and its strict chronological relation is
$x\prec y$ if and only if $x_1<y_1$ and $x_2<y_2$.
For the fixed anchors $a=(0,0)$ and $b=(1/2,1/2)$, let
\[
 A=\{x:a\prec x\prec b\}=(0,1/2)^2,\qquad
 v_\rho=\mu_\rho(A).
\]
Then $1/8\le v_\rho\le3/8$.
\end{proposition}
\begin{proof}
Integrating either marginal identity gives total mass one. The metric
matrix has determinant $-\rho^2$, so its volume density is $\rho$.
A nonzero future timelike tangent has both components positive, since
$g_\rho(\dot x,\dot x)=-2\rho\dot u\dot v<0$ and the time orientation
selects the positive branch. A future timelike curve therefore increases
both coordinates. Conversely, the straight segment between two points
with strictly increasing coordinates lies in the convex diamond and is
future timelike. This proves the relation and the expression for $A$.
Its ordinary area is $1/4$; integrating the density bounds gives the
asserted range. Boundaries have zero measure.
\end{proof}

Our units normalize the total geometry volume to one. Thus $v_\rho$ is a
dimensionless physical volume fraction. Without external scale
information this experiment cannot attach an absolute spacetime volume
unit to the answer.

Let the detector $\pi:\R^2\to\R$ be measurable, with the global bounds
\begin{equation}\label{eq:detector}
 1/R\le\pi(x)\le1,\qquad R\ge1.
\end{equation}
Here $R$ is known and deterministic. Write
\[
 Z=\int\pi\dd\mu_\rho>0,\qquad
 \nu(B)=Z^{-1}\int_B\pi\dd\mu_\rho,\qquad \theta=\nu(A).
\]
In the fixed-retained-count experiment, $X_1,\ldots,X_n$ are independent
with law $\nu$. The observation $Y_n$ consists of every directed strict
chronological relation among these events and the two anchors, modulo
permutations of sample labels fixing the anchor roles. This is a finite
measurable observation space. It includes sample--sample, anchor--sample
and anchor--anchor relations. Coordinates, rejected locations, detector
values, generated counts and clocks are not observed in this experiment.

The count
\[
 K=\sum_{i=1}^n\ind\{a\prec X_i\prec b\}
\]
is determined by $Y_n$ and unchanged by relabeling. Under the model,
$K\sim\Bin(n,\theta)$. The complete observation can contain more
information than $K$; only the upper procedure discards that information.
Write $P_{\rho,\pi}^{(n)}$ for the law of $Y_n$.

\begin{theorem}[Uniform joint accuracy]\label{thm:main}
For $n\ge1$ and $1\le R\le2$, put
\[
 a_n=n^{-1/2},\qquad b_R=\frac{R-1}{R+1},\qquad
 s(n,R)=\min\{1,a_n+b_R\}.
\]
For every admitted $(\rho,\pi)$, the observable fraction $\widehat v=K/n$
satisfies
\begin{equation}\label{eq:upper}
 P_{\rho,\pi}^{(n)}\{|\widehat v-v_\rho|>2s(n,R)\}\le1/20.
\end{equation}
Conversely, let $(\Omega,\xi)$ be any probability space and let
$T:Y_n\times\Omega\to\R$ be any measurable estimator. There exists an
admitted $(\rho,\pi)$ such that
\begin{equation}\label{eq:lower}
 (P_{\rho,\pi}^{(n)}\otimes\xi)
 \{|T-v_\rho|>s(n,R)/256\}\ge1/4.
\end{equation}
The seed law $\xi$ is common to the alternatives and independent of the
observation. No boundedness or determinism restriction is placed on $T$.
\end{theorem}

The proof appears in Section~\ref{sec:rate}. If the 95\% minimax radius is
the infimum over eligible procedures of radii with worst-case failure at
most $1/20$, this theorem places it between $s(n,R)/256$ and $2s(n,R)$.
It holds at $R=1$ and uniformly for sequences $R=R_n$ tending to one.
The constants are deliberately conservative and are not asserted optimal.
Section~\ref{sec:report} supplies the more useful reporting procedure.
Zero samples are treated separately in Proposition~\ref{prop:empty}.

\section{What independent thinning produces}\label{sec:thinning}
We give the process argument on an arbitrary probability space
$(E,\mu)$ with a measurable $0\le\pi\le1$ and $Z=\int\pi\dd\mu>0$.
Generate independent pairs $(X_i,U_i)$, where $X_i\sim\mu$ and
$U_i\sim\operatorname{Unif}[0,1]$, with all coordinates independent.
Retain event $i$ precisely when $U_i<\pi(X_i)$, and list retained
locations in increasing \emph{generation index}. This order is not a
sorting by spacetime location.

\begin{proposition}[Finite and mixed generation]\label{prop:finite-thinning}
Let $\alpha(B)=\int_B\pi\dd\mu$ and $\nu=\alpha/Z$. For fixed generated
count $N$, a retained index set $I=\{i_1<\cdots<i_n\}$, and measurable
$B_1,\ldots,B_n$, the joint probability of exactly this index set and
$X_{i_j}\in B_j$ for every $j$ is
\begin{equation}\label{eq:pattern}
 Z^n(1-Z)^{N-n}\prod_{j=1}^n\nu(B_j).
\end{equation}
Writing $C_N$ for the retained count and $V_n$ for its ordered locations,
their joint subprobability measure is
\begin{equation}\label{eq:counttuple}
 \Prb(C_N=n,V_n\in\cdot)
 =\binom Nn Z^n(1-Z)^{N-n}\nu^{\otimes n}(\cdot)
 \quad(0\le n\le N).
\end{equation}
If $N$ is chosen independently with any probability mass function $q_N$,
the right side is summed against $q_N$. Conditional on any retained count
of positive probability, the locations consequently have law
$\nu^{\otimes n}$.
\end{proposition}
\begin{proof}
Tonelli's theorem gives
\[
 \Prb(X_i\in B,\ U_i<\pi(X_i))
 =\int_B\int_0^1\ind\{u<\pi(x)\}\dd u\dd\mu(x)=\alpha(B).
\]
The rejection probability is $1-Z$. Multiplying these one-event
identities over the independent generated pairs proves
\eqref{eq:pattern}. There are $\binom Nn$ disjoint retained patterns;
the same product law for the ordered locations occurs for each. Summing
proves \eqref{eq:counttuple} on measurable rectangles. Finite measures
agreeing on such rectangles agree on the product sigma algebra, so the
statement holds for every measurable tuple event. Independence of $N$
then gives the mixture by the law of total probability. Division by its
count mass is justified only when that mass is positive.
\end{proof}

These are submeasure identities even for impossible patterns. In
particular, they include $Z=1$ and $n=0$, using empty products and $0^0=1$.
The empty tuple has its unique probability law; impossible counts are not
assigned a conditional distribution by division by zero. An arbitrarily
location-dependent choice of $N$ is outside the independence assumption.

\begin{proposition}[Poisson generation and its Laplace functional]
\label{prop:poisson}
If $N\sim\Poi(\kappa)$ independently, with $\kappa\ge0$, then
$C_N\sim\Poi(\kappa Z)$ and, conditional on a possible count $n$, the
retained tuple has law $\nu^{\otimes n}$. For every nonnegative measurable
real function $h$,
\begin{equation}\label{eq:laplace}
 \E\exp\left(-\sum_{i\text{ retained}}h(X_i)\right)
 =\exp\left(\kappa\int\pi(x)(e^{-h(x)}-1)\dd\mu(x)\right).
\end{equation}
\end{proposition}
\begin{proof}
For fixed $n$, set $N=n+r$ in the mixture of \eqref{eq:counttuple}.
The scalar mass is
\[
 \sum_{r=0}^\infty e^{-\kappa}\frac{\kappa^{n+r}}{(n+r)!}
 \binom{n+r}{n}Z^n(1-Z)^r
 =e^{-\kappa Z}\frac{(\kappa Z)^n}{n!}.
\]
The tuple factor does not depend on $r$. This proves the count and
conditional laws, also at $\kappa=0$. Alternatively a single generated
pair contributes a Laplace factor
\[
 M=\int(1-\pi+\pi e^{-h})\dd\mu
   =1+\int\pi(e^{-h}-1)\dd\mu\in[0,1].
\]
For $N$ pairs the expectation is $M^N$. Summing its Poisson series yields
$e^{\kappa(M-1)}$, which is \eqref{eq:laplace}. All integrands are
bounded or nonnegative, so the integrations and series are justified.
\end{proof}

Repeated locations, if the underlying measure has atoms, are counted with
multiplicity. For a nonnormalized geometry of total volume $V$, taking
$\mu=\mathrm{vol}/V$ and generated intensity $\kappa V$ instead produces
the retained intensity measure $\kappa\pi\,\mathrm{vol}$. This identity
does not make $V$ identifiable from a fixed retained order.

\begin{proposition}[Stopping after the first $n$ retained events]
\label{prop:stopping}
In the infinite independent stream, the prefix ending at the $n$th retained
event is almost surely finite for every $n\ge1$. Its ordered retained
tuple has law $\nu^{\otimes n}$. For $n=0$, the prefix and tuple are empty.
\end{proposition}
\begin{proof}
After any fixed index, the probability that the next $m$ trials all fail is
$(1-Z)^m$, which tends to zero. By continuity of probability on decreasing
events, having no later success has probability zero. If there were only
finitely many successes, some integer index would be after the last one.
A countable union of these null events is null. Thus every required
success occurs in a finite prefix.

For $n\ge1$ and a terminal prefix of length $N\ge n$, its last event is
accepted and exactly $n-1$ of the preceding $N-1$ events are accepted.
Applying \eqref{eq:pattern} to these disjoint patterns gives
\[
 \Prb(T_n=N,V_n\in B)
 =\binom{N-1}{n-1}Z^n(1-Z)^{N-n}\nu^{\otimes n}(B).
\]
These terminal events are disjoint and exhaust probability one by the
first paragraph. Summing first with $B$ the whole tuple space shows their
scalar masses sum to one, and then summing for arbitrary $B$ proves the
result. At $n=0$ the assertion is the empty-product identity.
\end{proof}

The accepted formal proof sums a finite set of terminal patterns instead
of simplifying its cardinality to $\binom{N-1}{n-1}$. The displayed
combinatorial simplification is elementary prose packaging; the actual
stopped-tuple law is formally verified. The stopping argument concerns a
prescribed retained count, not a general data-dependent stopping rule.

\section{A finite report for physical volume}\label{sec:report}
The following transformation separates detector uncertainty from sampling
uncertainty. Here $R\ge1$ may exceed the range needed in the joint rate.
Define, for $x\in[0,1]$,
\begin{equation}\label{eq:transforms}
 L_R(x)=\frac{x}{R-(R-1)x},\qquad
 U_R(x)=\frac{Rx}{1+(R-1)x}.
\end{equation}

\begin{proposition}[Bounded selection]\label{prop:selection}
For a probability measure $\mu$, a measurable target set of probability
$v$, and a detector $\pi\in[1/R,1]$, let $\theta$ be its retained
probability. Then
\begin{equation}\label{eq:identification}
 L_R(\theta)\le v\le U_R(\theta),\qquad |\theta-v|\le b_R.
\end{equation}
Both transforms are increasing on $[0,1]$.
\end{proposition}
\begin{proof}
Put $s=\int_A\pi\dd\mu$ and $t=\int_{A^c}\pi\dd\mu$. Then
$v/R\le s\le v$ and $(1-v)/R\le t\le1-v$.
Since $s+t>0$ and $\theta=s/(s+t)$, these inequalities imply
\[
 \frac{v}{v+R(1-v)}\le\theta\le
 \frac{Rv}{Rv+1-v}.
\]
Solving the two inequalities for $v$ proves the bounds in
\eqref{eq:identification}, including $v=0,1$. All relevant denominators
are positive. For $x\le y$, the differences of the transforms are
\[
 L_R(y)-L_R(x)=\frac{R(y-x)}{[R-(R-1)x][R-(R-1)y]},
\]
\[
 U_R(y)-U_R(x)=\frac{R(y-x)}{[1+(R-1)x][1+(R-1)y]},
\]
which prove monotonicity. Finally, direct common-denominator calculations
give
\begin{align*}
 b_R-(x-L_R(x))&=
 \frac{(R-1)[R(1-x)^2+x^2]}{(R+1)[R-(R-1)x]}\ge0,\\
 b_R-(U_R(x)-x)&=
 \frac{(R-1)[(1-x)^2+Rx^2]}{(R+1)[1+(R-1)x]}\ge0.
\end{align*}
Substitute $x=\theta$ and use the established bounds for $v$.
\end{proof}

These binary transforms appear explicitly in Aronow and Lee
\cite[p.~237]{aronow}. For observed zero--one values with $k$ ones, a
normalized weighted mean with weights in $[1,R]$ has minimum $L_R(k/n)$
and maximum $U_R(k/n)$: assign the smallest weights to the ones and the
largest to the zeros for the minimum, and reverse them for the maximum.
The analogous probability bounds are sharp when no further restrictions
are imposed, by constant detectors on $A$ and $A^c$. Neither fact proves
sharpness of the identified set inside $\K$ conditional on the full order.
Moreover, plug-in bounds alone are not a confidence interval. With
$n=1$, $R=1$ and the flat density, the plug-in interval is $\{0\}$ or
$\{1\}$, and has zero coverage for the true value $1/4$.

\begin{theorem}[Finite confidence from valid endpoint tables]
\label{thm:report}
Let $0<\delta<1$ and choose endpoints $0\le l(k)\le u(k)\le1$ for
every $k\in\{0,\ldots,n\}$. Suppose, with $X_p\sim\Bin(n,p)$, that
\begin{align}
 l(k)=0\quad&\text{or}\quad\Prb(X_{l(k)}\ge k)\le\delta/2,\label{eq:tailchecks}\\
 u(k)=1\quad&\text{or}\quad\Prb(X_{u(k)}\le k)\le\delta/2.\nonumber
\end{align}
If $K\sim\Bin(n,\theta)$, then
\[
 \Prb\{l(K)\le\theta\le u(K)\}\ge1-\delta.
\]
In the geometric model, a physical-volume report with the same coverage is
\begin{equation}\label{eq:physical-report}
 I(K)=[L_R(l(K)),U_R(u(K))]\cap[1/8,3/8].
\end{equation}
At $n=0$, take $[l(0),u(0)]=[0,1]$, giving $I(0)=[1/8,3/8]$.
The coverage holds in each conditional retained-count experiment of
Section~\ref{sec:thinning}, in the first-$n$ experiment, and after mixing
over retained counts when the same failure budget is used at each count.
\end{theorem}
\begin{proof}
Couple $X_p=\sum_{i=1}^n\ind\{U_i<p\}$ with independent uniform
$U_i$. This count increases with $p$, so its upper tails increase and
its lower tails decrease with $p$. Under the true $\theta$, a failure
$l(k)>\theta$ implies
$\Prb(X_\theta\ge k)\le\delta/2$ by \eqref{eq:tailchecks} and
monotonicity. For any integer-valued random variable $X$, the event
$\{\Prb(X\ge X_{\rm obs})\le c\}$ has probability at most $c$:
if there is a qualifying integer in its finite support, let $k_0$ be the
smallest; all qualifying outcomes are in $\{X\ge k_0\}$, whose mass is
at most $c$. If none qualifies the event is empty. Applying this with
$c=\delta/2$ bounds the lower-endpoint failure. The largest qualifying
integer and a lower tail give the same bound for $u(k)<\theta$.
The union bound proves the first assertion. No monotonicity of the
endpoint table as a function of $k$ is required.

On the successful binomial event, Proposition~\ref{prop:selection} and
monotonicity imply $L_R(l(K))\le v_\rho\le U_R(u(K))$.
The deterministic range in Proposition~\ref{prop:geometry} permits
intersection without losing that event. An empty intersection stays empty;
it has no target coverage. The zero-sample statement is immediate.
The tuple laws already proved transfer this argument to each conditional
or stopped experiment. Finally, a mixture of failure probabilities each
at most $\delta$ is at most $\delta$.
\end{proof}

\begin{proposition}[Integer endpoint certificates]\label{prop:certificate}
For $p=m/d$ with integers $0\le m\le d$ and $d>0$, binomial probabilities
can be certified exactly using
\[
 \Prb(X_p=j)=\frac{t_j}{d^n},\qquad
 t_j=\binom njm^j(d-m)^{n-j}.
\]
For $0<m<d$ and $j<n$, the recurrence
\begin{equation}\label{eq:recurrence}
 t_{j+1}=\frac{t_j(n-j)m}{(j+1)(d-m)}
\end{equation}
has exact integer division. Rational tail comparisons in
\eqref{eq:tailchecks} therefore reduce to integer comparisons.
\end{proposition}
\begin{proof}
Substitution in the binomial mass formula gives the common denominator
$d^n$. The identity $(j+1)\binom n{j+1}=(n-j)\binom nj$ proves
\eqref{eq:recurrence} by multiplication; its denominator is positive and
the quotient equals the integer $t_{j+1}$. Sum these numerators over the
desired tail and cross-multiply positive denominators against the rational
budget. At $m=0$ or $m=d$, use the deterministic count at $0$ or $n$
instead of dividing by zero. The empty-sample mass is one.
\end{proof}

The implementation proposes Clopper--Pearson endpoints using floating
point beta quantiles, rounds them outward to denominator $2^{16}$, then
checks the rational tail inequalities exactly. Proposal precision is not
the certificate: a failed check triggers an outward adjustment or a
full-range fallback. Both preserve the stated coverage when the final
table passes. The recorded pilot checks every count in its tables; its
adjacent inward checks additionally bound the outward grid error for those
tables. This is not a guarantee on every possible proposal/fallback path.

For comparison we also use
$[\max(0,K/n-r_n),\min(1,K/n+r_n)]$ with a rational
$r_n\ge\sqrt{2/n}$. Hoeffding's bound is at most $2e^{-4}<1/20$:
indeed $e^4\ge\sum_{j=0}^5 4^j/j!>40$. Transforming these endpoints
gives a valid analytic physical report. The implementation chooses the
rational radius by an integer square-root calculation. At $n=0$ it uses
the no-data interval. Python execution and floating point proposal
semantics are tested software, not compiled Lean theorems.

\section{Proof of the joint accuracy theorem}\label{sec:rate}
We first give the upper argument, then prove two lower obstructions for
the actual marked observation.

\subsection{Uniform upper bound}
By Proposition~\ref{prop:selection}, $|\theta-v_\rho|\le b_R$.
The count fraction satisfies $0\le K/n\le1$ and $v_\rho\in[0,1]$.
If $a_n+b_R\ge1$, error strictly greater than $2s=2$ is impossible.
Otherwise, error greater than $2(a_n+b_R)$ forces
$|K/n-\theta|>2a_n+b_R\ge2a_n$. Hoeffding's inequality therefore gives
\[
 \Prb\{|K/n-v_\rho|>2s(n,R)\}
 \le2\exp(-2n(2a_n)^2)=2e^{-8}\le1/20.
\]
This proves \eqref{eq:upper}; the same argument works for any $R\ge1$.
The restriction $R\le2$ is used below in the geometric detector alternative.

\subsection{Testing, geometry and actual likelihoods}
\begin{lemma}[Testing with an independent seed]\label{lem:testing}
Let $P,Q$ be probability laws on a finite observation space, and suppose
$\TV(P,Q)\le\eta$. For any common independent seed law $\xi$, measurable
real estimator $T$ and targets $v_0<v_1$, if $2r<v_1-v_0$, then
\[
 \max\{(P\otimes\xi)(|T-v_0|>r),
          (Q\otimes\xi)(|T-v_1|>r)\}\ge(1-\eta)/2.
\]
In particular, identical laws give the lower bound $1/2$.
\end{lemma}
\begin{proof}
Let $B_i=\{|T-v_i|\le r\}$. These measurable sets are disjoint:
membership in both would give $v_1-v_0\le2r$. For any measurable set
$B$ in the observation--seed product, its section probabilities
$h(y)=\xi\{\omega:(y,\omega)\in B\}$ lie in $[0,1]$. The finite
signed-mass formula gives
\[
 |(P\otimes\xi)(B)-(Q\otimes\xi)(B)|
 =\left|\sum_y(P\{y\}-Q\{y\})h(y)\right|\le\TV(P,Q).
\]
Thus $(P\otimes\xi)(B_0)+(Q\otimes\xi)(B_1)
\le(Q\otimes\xi)(B_0\cup B_1)+\eta\le1+\eta$.
The complementary failure probabilities sum to at least $1-\eta$,
which proves the claim. The identical-law argument also works on an
arbitrary measurable observation space.
\end{proof}

\begin{lemma}[An admitted polynomial family]\label{lem:family}
For $0\le\varepsilon\le1/2$, define
\[
 f(t)=2t-1,\qquad
 \rho_\varepsilon(u,v)=1+\varepsilon f(u)f(v).
\]
Then $\rho_\varepsilon\in\K$ and
\begin{equation}\label{eq:family-target}
 v_{\rho_\varepsilon}=1/4+\varepsilon/16.
\end{equation}
Under $n$ unthinned flat coordinate samples, the likelihood of the
$\rho_\varepsilon$ sample is $L=\prod_i\rho_\varepsilon(X_i)$, with
\begin{equation}\label{eq:likelihood}
 \E_0L=1,\qquad \E_0(L-1)^2=(1+\varepsilon^2/9)^n-1.
\end{equation}
\end{lemma}
\begin{proof}
The polynomial is smooth. On $[0,1]$, $|f|\le1$, giving
$1-\varepsilon\le\rho_\varepsilon\le1+\varepsilon$.
For $x=(u,v)$ and $y=(u',v')$ in $D$, adding and subtracting
$f(u')f(v)$ yields
\[
 |\rho_\varepsilon(x)-\rho_\varepsilon(y)|
 \le2\varepsilon(|u-u'|+|v-v'|)
 \le4\varepsilon\|x-y\|_2\le2\|x-y\|_2.
\]
Since $\int_0^1 f=0$, both coordinate marginals equal one.
Also $\int_0^{1/2}f=-1/4$, proving \eqref{eq:family-target}.
The flat law is uniform on $D$. Multiplication of the single-point
density identities gives the actual product likelihood on coordinate
space; it is nonnegative on that support and integrates to one.
Because $\int_0^1f^2=1/3$, the single-point second moment is
$1+\varepsilon^2/9$. Independence gives its $n$th power for $\E_0L^2$.
Subtracting $2\E_0L-1=1$ proves \eqref{eq:likelihood}.
\end{proof}

\begin{lemma}[Sampling obstruction]\label{lem:sampling}
For every $n\ge1$, $R\ge1$ and estimator/seed in
Theorem~\ref{thm:main}, there is an admitted model with unit detector
such that
\[
 (P_{\rho,1}^{(n)}\otimes\xi)(|T-v_\rho|>a_n/128)\ge3/8.
\]
\end{lemma}
\begin{proof}
Compare $\rho_0=1$ with $\rho_\varepsilon$ for
$\varepsilon=1/(2\sqrt n)\le1/2$, using $\pi=1$ for both.
This detector is allowed for every $R\ge1$.
Put $x=1/(36n)$. Bernoulli's inequality gives
$(1-x)^n\ge1-nx=35/36$; it follows by induction from
$(1-x)^{j+1}\ge(1-jx)(1-x)\ge1-(j+1)x$ whenever needed here.
Moreover $(1+x)^n(1-x)^n=(1-x^2)^n\le1$. Hence
\[
 \E_0(L-1)^2=(1+1/(36n))^n-1\le1/35.
\]
For coordinate laws with density $L$ relative to the flat law, their
total variation is $\tfrac12\E_0|L-1|$; this follows by splitting the
zero-integral function $L-1$ into positive and negative parts.
Cauchy--Schwarz therefore bounds it by $1/(2\sqrt{35})\le1/4$.

Mapping coordinates to the full marked order cannot increase total
variation: the inverse image of an observation event is a coordinate
event. The two full marked-order laws thus have total variation at most
$1/4$, regardless of the other relations present in the observation.
Their targets differ by $\varepsilon/16=a_n/32$. With $r=a_n/128$,
the strict separation $2r<a_n/32$ holds. Lemma~\ref{lem:testing} now
gives the asserted lower bound $(1-1/4)/2=3/8$.
\end{proof}

\subsection{Exact detector confounding}
\begin{lemma}[Calibration obstruction]\label{lem:detector}
For $1<R\le2$, $n\ge0$ and any eligible estimator/seed, there is an
admitted model satisfying
\[
 (P_{\rho,\pi}^{(n)}\otimes\xi)(|T-v_\rho|>b_R/64)\ge1/2.
\]
The two alternatives used in the proof have the same unnormalized
accepted-location measure. Their finite generated-count and stopped
accepted-tuple experiments also have the same laws.
\end{lemma}
\begin{proof}
Set $\varepsilon=b_R\in(0,1/3]$. Compare the flat density with
$\rho_\varepsilon$ from Lemma~\ref{lem:family}. Define the detectors
globally by
\[
 \pi_0(x)=1-\varepsilon,\qquad
 \pi_\varepsilon(x)=
 \frac{1-\varepsilon}
 {\max\{1-\varepsilon,\min\{1+\varepsilon,\rho_\varepsilon(x)\}\}}.
\]
The denominator is continuous and at least $1-\varepsilon>0$. Thus
$\pi_\varepsilon$ is measurable and everywhere lies between
$(1-\varepsilon)/(1+\varepsilon)=1/R$ and $1$.
Also $1-\varepsilon\ge1/R$, so $\pi_0$ is admitted. The clamping
matters outside $D$; an unclamped reciprocal need not be a valid global
detector.

On $D$ the denominator equals $\rho_\varepsilon$, so
$\rho_\varepsilon\pi_\varepsilon=1-\varepsilon
=\rho_0\pi_0$. Both unnormalized accepted measures equal
$(1-\varepsilon)\,\mathrm{Leb}|_D$, both acceptance masses equal
$1-\varepsilon$, and both retained laws are uniform on $D$.
Their full marked-order laws are therefore identical for every $n$.
Their target volumes differ by $\varepsilon/16$, which is strictly
larger than $2(\varepsilon/64)$. The identical-law case of
Lemma~\ref{lem:testing} proves the claim.

Equality of the accepted submeasures gives equality of $Z$ and $\nu$.
Equations~\eqref{eq:pattern}--\eqref{eq:counttuple} then give equality
even if the generated count and acceptance pattern are retained, provided
the same generated-count law is used. The mixed, Poisson and stopped
identities give their respective equalities as well. Rejected locations
are not part of these equal-law records.
\end{proof}

\subsection{Combining the obstructions and treating zero samples}
\begin{proof}[Completion of Theorem~\ref{thm:main}]
Only the lower bound remains. If $a_n\ge b_R$, then
$s(n,R)\le2a_n$ and $s(n,R)/256\le a_n/128$.
Lemma~\ref{lem:sampling} supplies failure at least $3/8$, hence at
least $1/4$, at the smaller radius in \eqref{eq:lower}.
If $b_R>a_n$, then $R>1$, $s(n,R)\le2b_R$, and
$s(n,R)/256\le b_R/128<b_R/64$.
Lemma~\ref{lem:detector} supplies failure at least $1/2$ at the larger
radius, and therefore at least $1/4$ at the stated radius. These cases
exhaust all $n\ge1$ and $1\le R\le2$, including $R=1$.
\end{proof}

\begin{proposition}[Zero retained samples]\label{prop:empty}
For $n=0$, the constant estimator $1/4$ has error at most $1/8$ over
$\K$. For every $R\ge1$ and every measurable randomized estimator of
the empty-sample marked observation, some admitted model satisfies
\[
 (P_{\rho,\pi}^{(0)}\otimes\xi)(|T-v_\rho|>1/128)\ge1/2.
\]
\end{proposition}
\begin{proof}
The upper statement follows from $v_\rho\in[1/8,3/8]$.
For the lower statement, compare the flat density and
$\rho_{1/2}$, both with unit detector. There are no random sampled
points. The observation still includes the fixed anchor relations, but
these are identical in both models, so its law is the same Dirac mass.
The targets differ by $1/32>2/128$. Apply the identical-law testing
argument with $r=1/128$ and the arbitrary independent seed.
\end{proof}

\section{Two preserved computational studies}\label{sec:studies}
The following studies were frozen and executed before preparation of this
manuscript. Their data are preserved unchanged. They ask different
questions and do not form one pooled experiment. All reported widths are
dimensionless volume-fraction widths. An expected width averages over a
specified count law; a realized width is the width at one observed count.
Coverage means probability of containing the true target and is distinct
from either width.

\subsection{Original generic-probability pilot}
The pilot used 300 combinations of retained count, target probability,
detector ratio and detector pattern, and four methods, producing 1,200
rows. Counts were $n\in\{0,64,256,1024\}$. Each configuration used
2,000 Monte Carlo counts as well as a numerical sum over the complete
binomial law. At $n=0,64$, 600 method rows also received exact rational
coverage checks; 450 belonged to the three valid methods and 150 to the
deliberately misspecified method that ignored detection bias.
The generic target grid includes probabilities outside $[1/8,3/8]$;
it is not a grid confined to our fixed geometric interval.

The valid methods were the transformed certified binomial interval, the
transformed analytic interval, and the full $[0,1]$ no-data interval.
The fourth method estimated the detected probability as if it were the
physical probability. The 1,348 distinct count intervals passed their
exact endpoint checks, with 2,688 adjacent inward endpoint rejections and
no adjustments or fallbacks. Minimum complete-law numerical coverage
over valid methods was $0.951071$; the minimum exact rational coverage
in the checked small-count subset was $0.966807$.
The bias-ignoring method fell to approximately $1.20\times10^{-19}$.

\begin{table}[htbp]
\centering
\begin{tabular}{rrr}
\toprule
$R$ & Certified binomial & Analytic \\
\midrule
1 & 0.06214 & 0.08839 \\
1.1 & 0.10938 & 0.13539 \\
1.25 & 0.17195 & 0.19746 \\
1.5 & 0.25874 & 0.28320 \\
2 & 0.38717 & 0.40939 \\
\bottomrule
\end{tabular}
\caption{Original pilot: worst expected widths on its prescribed generic-probability grid at $n=1024$. The no-data width is one.}
\label{tab:pilot}
\end{table}

Table~\ref{tab:pilot} reports the worst \emph{expected} width over this
pilot's prescribed target and detector patterns at $n=1024$. Its frozen
width threshold of $0.20$ passed for all tested cells with $R\le1.25$.
This was a bounded feasibility result, not a claim over untested targets
or a field detector. Execution took 277.11 seconds, with 4,774,983 bytes
of Python-traced peak memory. That metric does not include all native
allocations or measure total process resident memory.

The pilot separately generated four actual geometric configurations, each
with 200 trials of $n=1024$ retained events: flat density or
$\rho_{1/4}$, with constant or compensating detection. Their targets
were respectively $1/4$ and $17/64$. Recorded coverage proportions were
$0.935$, $0.955$, $0.965$ and $1.000$ in the preserved artifact order.
The first is 187 successes out of 200; it has not been replaced or hidden
because it lies below $0.95$. Its exact binomial confidence interval for
the repeated-trial success probability is approximately $[0.8914,0.9650]$;
the complete-law coverage for the corresponding flat case is about
$0.952905$. Thus this Monte Carlo fluctuation does not contradict the
finite guarantee. The compensated polynomial case has $R=5/3$ and mean
width about $0.24188$.

\subsection{Frozen geometric finite comparison}
The later study used the known geometric range $[1/8,3/8]$ for every
method. It compared four reports from the same $(n,K,R)$: transformed
certified binomial, transformed Hoeffding, the joint-rate interval centered
at $K/n$ with radius $2s(n,R)$, and the fixed known range. Every interval
was intersected with the known range. An empty intersection had zero width
and no coverage; it was not replaced by a fallback. All methods returned
the known range at $n=0$.

The formula grid used $n=q^2$ for $q=1,2,4,8,16,32,64$ and
\[
 R\in\{1,1025/1024,65/64,17/16,9/8,5/4,3/2,2\},
\]
augmented by admitted crossover values with
$b_R\in\{1/(2q),1/q,\min(2/q,1/3)\}$ and $b_R\le1/3$.
There were 71 rate cells. Report calculations included every count for
$n=0,1,4,16,64$ and seven selected counts
$0,1,n/4,n/2,3n/4,n-1,n$ for $n=256,1024$.
All 104 distinct binomial endpoint reports passed exact tail checks with
no fallback.

For complete-law coverage, the study used
$\rho_\varepsilon$ with $\varepsilon\in\{0,1/4,1/2\}$ and uniform,
inside-low or inside-high piecewise detection in $[1/R,1]$.
For $n\le64$, all 1,656 exact rational coverage checks passed; their
minimum was approximately $0.96288634$. Expected widths were enclosed
by outward rational display intervals of width at most $10^{-12}$,
and empty-report probabilities were computed exactly. Complete-law
coverage was not calculated at $n=256$ or $1024$ in this study.

Across 1,108 paired realized-report cases, the joint-rate interval was
wider than the certified binomial report in 732 and equal in 376; it was
never narrower on that grid. Table~\ref{tab:finite} shows three selected
\emph{realized} widths at $n=1024$, $K=256$. This is evidence for retaining
the exact-binomial practical report, not a proof of universal interval
dominance or shortest possible intervals.
\begin{table}[htbp]
\centering
\begin{tabular}{rrrrr}
\toprule
$R$ & Certified binomial & Hoeffding & Joint rate & Known range \\
\midrule
1 & 0.053970 & 0.088388 & 0.125000 & 0.250000 \\
1.25 & 0.137225 & 0.170836 & 0.250000 & 0.250000 \\
2 & 0.249039 & 0.250000 & 0.250000 & 0.250000 \\
\bottomrule
\end{tabular}
\caption{Separate geometric study: realized widths at $n=1024$, $K=256$, after intersection with the known range.}
\label{tab:finite}
\end{table}


The formula checks included eight polynomial parameter values, product
divergence through $n=4096$, and five negative controls: an incorrect
target denominator, an incorrect second moment, an invalid strict-gap
boundary, an unclamped ambient detector and a shrunken binomial interval.
No Monte Carlo samples, numerical quadrature or random seeds were used in
this second study. Execution took 199.33 seconds with 59,031,222 bytes
of Python-traced peak memory. A separate artifact audit took about
3.15 seconds and checked the saved cells, including 261 binomial-mass
recurrences. The first attempt stopped because SciPy was unavailable;
its partial output remains preserved. The successful retry used the same
frozen protocol and runner after restoring the required dependency.
It did not alter the grid or select a favorable resampling outcome.

\section{Source comparisons and verification scope}\label{sec:verification}
Aronow and Lee's journal article and two-page supplement were read in
full, including the article's Proposition~1 and the supplement's Lemma~1
and Propositions~2--3. Their weighted-mean bounds, binary formulas and
asymptotic convergence arguments receive attribution here. They do not
themselves supply this paper's finite binomial confidence calibration or
the full marked-order geometric lower construction. In the reviewed
comparison, their binary transform is algebraically identical to the
transform already applied to the calibrated binomial benchmark. Adding
another numerical column with those same endpoints would duplicate the
existing method.

The related partial-identification comparisons use the editions actually
inspected: Imbens and Manski's 2003 CeMMAP working paper, Stoye's
17 December 2008 preprint, and Tudball and colleagues' fourth arXiv
version together with the journal and supplementary proof comparison.
Their confidence constructions and sample constraints have their own
regularity and asymptotic assumptions. We have not imported such
assumptions to replace our finite count argument, nor inferred that every
possible partial-identification method is covered by this comparison.
An uninspected journal edition is not claimed textually identical to its
reviewed preprint. Purchased source PDFs are kept private and are not
redistributed with our source package.

The accepted Lean development includes exact types and axiom reports for
797 exports in its original full audit. That total covers the entire
accepted dependency collection, not 797 new theorems of this manuscript.
The authoritative run is
\begin{center}\small
\texttt{space-volumerate-acceptance-20261010T023105Z-5d08cf}.
\end{center}
It ran on Google Cloud with Lean 4.30.0 and the pinned mathlib revision
recorded in the source archive. It completed 3,097 root build jobs with
zero warnings. Its reports use only the standard axioms
\texttt{propext}, \texttt{Classical.choice} and \texttt{Quot.sound}.
The 198 captured source files, dependency identities, raw logs, receipt
and task-owned instance shutdown evidence are retained. Proof commit:
\begin{center}\small
\texttt{354fa1f8dc8746328077041ef79d3621837f99e7}.
\end{center}

The accompanying \texttt{formal-map.json} associates each labeled
mathematical claim with the exact accepted exports and source hashes.
It also identifies elementary prose packaging, such as the determinant
calculation and terminal-pattern cardinality, rather than representing
every sentence as a separate kernel theorem. Principal entries are:
\begin{center}
\small
\begin{tabular}{ll}
\toprule
Manuscript result & Principal Lean module(s)\\
\midrule
Finite/mixed thinning & \texttt{ThinningCount}, \texttt{ThinningMixture}\\
Poisson and stopped laws & \texttt{ThinningPoisson}, \texttt{ThinningStoppedLaw}\\
Finite confidence and certificates & \texttt{BinomialCalibration}, \texttt{BinomialArithmetic}\\
Physical report transfer & \texttt{BinomialGeometry}, \texttt{ThinningMixedCoverage}\\
Geometric alternatives & \texttt{ConfoundingGeometry}, \texttt{ConfoundingMarked}\\
Sampling lower bound & \texttt{VolumeRateLikelihood}, \texttt{VolumeRateSampling}\\
Uniform upper/lower theorem & \texttt{VolumeRateUpper}, \texttt{VolumeRateJoint}\\
Zero samples & \texttt{VolumeRateEmpty}\\
\bottomrule
\end{tabular}
\end{center}

Preparation reuses this immutable acceptance; it does not rerun Lean or
change any accepted proof. The source verifier checks the full captured
archive against the original proof commit and the current mapped modules
against that archive. Later repository integration files can contain
additional projects without changing the original acceptance evidence.
The Python reports and mathematical source comparison are separately
audited evidence, not kernel verification of Python, SciPy, the operating
system or an actual detector. Formal correctness also does not establish
that the model is a correct description of a physical instrument.

The reproducibility package supplies the manuscript source and build
instructions, exact formal map, original proof snapshot and acceptance
logs, table extraction, and frozen computational evidence. Its manifest
records file hashes. The manuscript is licensed CC BY 4.0; source-code
licenses and the original repository license are preserved. This is an
unpublished preparation with no newly assigned manuscript DOI.

\section{Implications and limits}\label{sec:implications}
With perfectly calibrated uniform detection ($R=1$), the theorem gives
the familiar $n^{-1/2}$ scale for this scalar target. For fixed $R>1$,
the calibration term remains positive as $n$ grows. Within the stated
model, using more marked-order relations cannot force the worst-case
volume error to vanish: the compensated geometries have exactly the same
retained law. This is an information limit of the observation experiment,
not a failure of one estimator.

The sampling and calibration terms become comparable at
$n^{-1/2}\asymp(R-1)/(R+1)$. To keep a sampling-order accuracy as $n$
increases, the ratio uncertainty must shrink at that order as well.
This gives a statistical way to distinguish the benefit of more retained
events from the benefit of better detector knowledge. It does not choose
a monetary optimum, specify a calibration instrument, or measure a
sensor's achievable $R$. If $R$ were inferred from finite calibration
data, its uncertainty would need an additional coverage argument and
failure-budget allocation.

Physical use would require evidence for the causal observation map,
the sampling and independent detection assumptions, a justified bound on
the detector ratio, stable anchor identification and any conversion from
normalized to absolute units. None is demonstrated by the synthetic
studies. Dependence between detections, geometry-dependent stopping,
measurement error in causal relations or sample-selected anchors are
outside the present theorem. The results do not establish reconstruction
in higher dimensions, curvature estimates, operational resource gains or
a new quantum inference method.

The practical outcome is an inspectable finite report together with a
proved explanation of its calibration limit. Open-source, decentralized
review and downstream use can test and improve these results; anticipated
uptake is not treated as validation. New applications should first show
that their observation and calibration assumptions fit the experiment,
then compare the report with an established method using the same data
and coverage requirements.

\begin{thebibliography}{9}
\raggedright
\interlinepenalty=10000
\bibitem{aronow}
P.~M. Aronow and D.~K. Lee.
Interval estimation of population means under unknown but bounded
probabilities of sample selection.
\emph{Biometrika} 100(1), 235--240 (2013).
\href{https://doi.org/10.1093/biomet/ass064}{doi:10.1093/biomet/ass064}.
Journal article and two-page supplementary proofs inspected.

\bibitem{clopper}
C.~J. Clopper and E.~S. Pearson.
The use of confidence or fiducial limits illustrated in the case of the
binomial. \emph{Biometrika} 26(4), 404--413 (1934).
\href{https://doi.org/10.1093/biomet/26.4.404}{doi:10.1093/biomet/26.4.404}.

\bibitem{hoeffding}
W. Hoeffding.
Probability inequalities for sums of bounded random variables.
\emph{Journal of the American Statistical Association} 58(301), 13--30 (1963).
\href{https://doi.org/10.1080/01621459.1963.10500830}{doi:10.1080/01621459.1963.10500830}.

\bibitem{imbens}
G.~W. Imbens and C.~F. Manski.
Confidence intervals for partially identified parameters.
CeMMAP Working Paper CWP09/03, 29 May 2003 (reviewed edition).
\url{https://cemmap.ac.uk/wp-content/uploads/2020/08/CWP0903.pdf}.
Related journal version: \emph{Econometrica} (2004),
\href{https://doi.org/10.1111/j.1468-0262.2004.00555.x}{doi:10.1111/j.1468-0262.2004.00555.x}.

\bibitem{stoye}
J. Stoye.
More on confidence intervals for partially identified parameters.
Author preprint, 17 December 2008 (reviewed edition).
\url{https://stoye.economics.cornell.edu/docs/stoye-ecma2009-pre.pdf}.
Related journal version: \emph{Econometrica} (2009),
\href{https://doi.org/10.3982/ECTA7347}{doi:10.3982/ECTA7347}.

\bibitem{tudball}
M.~J. Tudball, R.~A. Hughes, K. Tilling, J. Bowden and Q. Zhao.
Sample-constrained partial identification with application to selection
bias. \emph{Biometrika} 110(2) (2023).
\href{https://doi.org/10.1093/biomet/asac042}{doi:10.1093/biomet/asac042}.
Reviewed preprint: \href{https://arxiv.org/abs/1906.10159v4}{arXiv:1906.10159v4},
4 October 2021; journal and supplementary proof comparison recorded in
the companion source audit.
\end{thebibliography}
\end{document}
